\chapter{Discussion}
\label{ch:discussion}

\section{Why the Proteomic Analysis Found No Signal}
\label{sec:discussion-phase1a}

Neither the univariate screen nor any of the classifiers recovered an association between the
model-derived plaque phenotype and the plasma proteome. Two explanations are available, they
operate at different levels, and the design of this study cannot separate them.

\subsection{A Biological Explanation}

The first explanation is that established subclinical carotid plaque simply does not produce a
detectable systemic protein signature. Atherosclerosis develops over decades, and the lesions
imaged here are, by the nature of a population imaging cohort, clinically silent. There is no
acute event to detect: no rupture, no thrombus, no inflammatory flare of the kind that would
plausibly register in circulation. What is present instead is a slow, focal accumulation in a
vessel wall, small in absolute terms against the volume of plasma into which anything it releases
would be diluted. A process may be consequential without being measurable, and subclinical plaque
is a plausible instance of that: large in its eventual clinical implications, small in its
contemporaneous biochemical footprint.

The strongest support for this reading comes from within the analysis itself. Age produced a
visible gradient in the leading principal components of the same protein matrix and remained
significantly associated after correction for multiple testing, while plaque status produced
neither. The same cohort, the same measurements and the same procedure detect one and not the
other. The proteome, on this evidence, reflects systemic physiological state rather than the
presence of a focal lesion.

One consideration qualifies the dilution argument. Carotid plaque is generally regarded not as an
isolated local finding but as a marker of systemic atherosclerosis, since subclinical disease is
typically present in several vascular territories at once \cite{fernandezfriera2015pesa}; a
participant with carotid plaque is therefore likely to carry coronary and aortic disease as well. The relevant burden is therefore
larger than the imaged lesion, and the argument from lesion volume holds only to the extent that
the carotid finding is a poor index of total atherosclerotic burden. Subclinical atherosclerosis
at the carotid and femoral sites is associated with coronary calcification
\cite{laclaustra2016femoral}, which suggests the carotid finding indexes more than itself.

\subsection{A Complication from the Same Cohort}

A strong form of the biological explanation is difficult to sustain. In the same UK Biobank
cohort, \textcite{omarov2024deep} report a genome-wide association study of carotid plaque whose
downstream analyses implicate lipoprotein(a) and interleukin-6 signalling, both of which have
been pursued as drug targets. Both are measurable in plasma, and both are present in the
Olink panel used here, as are interleukin-6 receptor, apolipoprotein B, PCSK9 and several matrix
metalloproteinases. The proteins that a causal analysis identifies in this cohort were therefore
measured in this analysis, and none of them emerged.

The resolution lies in the direction of inference. Mendelian randomisation identifies causes, not
consequences: lipoprotein(a) and interleukin-6 are upstream determinants that drive plaque
formation over decades, and a cross-sectional proteomic association would recover them only if
their effect on prevalent plaque status were large enough to survive both attenuation and
correction across 2{,}923 tests.

Two clinical results reported while this thesis was being written illustrate the same gap from
the other direction, and they concern precisely these two pathways. ZEUS, a phase III trial of
the interleukin-6 antibody ziltivekimab in more than 6{,}300 patients with atherosclerotic
disease, chronic kidney disease and raised C-reactive protein
\cite{ridker2026zeusdesign}, achieved the expected reduction in interleukin-6 signalling but
reported a hazard ratio of 0.99 (95\,\% CI 0.88--1.11) for major adverse cardiovascular events
\cite{novonordisk2026zeus, yao2026zeuscommentary}. Lp(a)HORIZON, a phase III trial of
pelacarsen in 8{,}323 patients with elevated lipoprotein(a) and established cardiovascular
disease \cite{cho2025horizondesign}, lowered lipoprotein(a) but did not meet its primary
endpoint of reducing cardiovascular events \cite{novartis2026pelacarsen}. Both targets were engaged; neither produced the
anticipated clinical benefit.

Neither result overturns the genetic evidence, and neither should be read as doing so. Both
trials lowered a long-standing exposure late, in patients whose disease was established and who
were already receiving guideline-directed therapy. The trial designs are published, but the
outcome data rest on company announcements: neither trial had reported in full at the time of
writing, and both should be checked against the eventual publications. What they do is separate two claims that the
genetic evidence alone does not distinguish: that a molecule contributes to how disease develops
across a lifetime, and that its concentration at a given moment is informative about the disease
that has developed. The same distinction governs the analysis performed here. A protein may shape
how a lesion forms over decades without its level at a single blood draw carrying information
about whether that lesion is now present. A molecule can be causal and still be a poor
cross-sectional marker, and the two properties are established by different kinds of evidence.

The claim that the present data support is consequently narrower than the one stated above. It is not that carotid atherosclerosis leaves no trace in blood, but
that established subclinical plaque produces no proteomic signature detectable \emph{as a
consequence} of its presence, under the sample size and phenotype quality available here.

There is a further reason to expect a binary label to be the wrong instrument for this
particular question. Interleukin-6 concentration is associated not with the presence of carotid
plaque but with its echogenicity \cite{yamagami2004il6} -- that is, with a continuous property
of the lesion's composition rather than with whether one is there at all. A phenotype that
records only presence discards exactly that information. If the proteins implicated by the causal
analysis relate to what a plaque is made of rather than to whether it exists, then no amount of
statistical power applied to a binary label would recover them, and the appropriate phenotype
would be a quantitative description of the lesion -- which is what a segmentation model, unlike a
detector, is able to produce.

The nominally significant proteins are of limited help in adjudicating this. Several are
biologically plausible in context -- matrix metalloproteinase 12 is a macrophage elastase with an
established role in plaque remodelling, and latent TGF-$\beta$ binding protein 2 and Fas ligand
have connections to matrix turnover and apoptosis respectively. But all nine of the leading hits
have between 1.6\,\% and 5.4\,\% missing values, against a panel median of 15.4\,\%. Better-measured
proteins carry more information and therefore attract smaller $p$-values irrespective of biology,
and median imputation compresses the variance of sparsely measured proteins and further reduces
their detectability. The ranking is thus confounded with measurement completeness, and the
apparent biological coherence of the top of the list should not be read as evidence.

\subsection{A Measurement Explanation}

The second explanation concerns the phenotype rather than the biology. The plaque labels are not
measurements but model outputs, and a binary exposure measured with error biases estimated
associations towards the null when the error is unrelated to the outcome
\cite{bross1954misclassification, copeland1977bias}. For a binary exposure the attenuation is
exact rather than approximate: the risk difference that such an analysis recovers is the true
risk difference multiplied by the sensitivity and specificity of the classifier minus one --
that is, by Youden's $J$, the same quantity used to evaluate the detection models in
Chapter~\ref{ch:results}. The measure that evaluates the model is therefore also the factor by
which it attenuates whatever the downstream analysis is looking for. On its own UK Biobank test
set the base model attains $J = 0.787$, so roughly a fifth of any true association is lost before
the analysis begins, and more than that on data unlike its development cohort. Combined with a Bonferroni threshold of
$1.7 \times 10^{-5}$ and an analysis cohort of 2{,}542 participants, only a comparatively large
effect could have been detected.

The two explanations are not alternatives so much as statements about different links in the same
chain, and both may hold simultaneously. What matters for the interpretation of this thesis is
that its design cannot distinguish them. A null result is compatible with an absent signal and
with a signal destroyed in measurement, and nothing in the analysis performed here separates the
two cases.

Phase 1B was, in retrospect, precisely the experiment that would have separated them. Relating an
independently derived cardiovascular risk estimate to the same proteomic data would have placed
the question on a phenotype not produced by the imaging model: an association there, alongside
its absence for plaque status, would have implicated the phenotype; its absence in both would
have implicated the proteomic data or the size of the cohort. That this comparison could not be
carried out is the reason the question remains open.

\subsection{An Alternative Design}

Independently of which explanation holds, the design tested prevalent plaque against
contemporaneously measured protein levels. If the proteome reflects the rate at which
atherosclerosis is progressing rather than the burden that has accumulated, a cross-sectional
comparison cannot detect it in principle. A design relating protein levels to incident plaque, or
to progression between two imaging visits, addresses a different and arguably more appropriate
question, and UK Biobank is in the process of assembling the data such a design would need: the
repeat imaging project aims to re-scan 60{,}000 participants, of whom roughly 20{,}000 had
attended by the end of 2025 \cite{ukbiobank_repeat_imaging}. Whether that resource supports a
plaque-progression analysis depends on the repeat protocol including carotid ultrasound, which
is a narrower question than whether repeat imaging exists.

\section{External Validation and Domain Shift}
\label{sec:discussion-domain-shift}

The base model loses 6.2 percentage points of sensitivity and 26.7 of specificity when moved from
the cohort it was developed on to Ar-PlaqSegm1. The asymmetry is the informative part. The model
still finds the plaques that are present; what it loses is the ability to recognise their absence.

\subsection{The Loss Is Not a Threshold Artefact}

An obvious objection is that the confidence threshold of \textcite{omarov2024deep} was tuned on
UK Biobank and simply does not transfer. The data answer this. At the threshold selected on the
validation split the base model reaches $J = 0.458$; the best value attainable anywhere on the
test set, with the threshold chosen on the test set itself, is 0.500. Re-tuning recovers 0.042 of
the 0.329 that separates external from in-cohort performance. The remainder is a change in what
the model represents, not in where its decision boundary is placed.

\subsection{Why Specificity Rather Than Sensitivity}

The composition of the external dataset offers a concrete explanation. Ar-PlaqSegm1 enrolled
patients with mild, moderate or severe atherosclerosis and a vascular wall thickness exceeding
1\,mm (Section~\ref{subsec:external-dataset}). Its plaque-free images are therefore not images of
healthy vessels: they are images of thickened, diseased vessels that happen not to contain a
focal plaque. In a population cohort the plaque-free class consists largely of unremarkable
vessels, and a model trained there learns a correspondingly permissive notion of what counts as
normal.

This is consistent with what the fine-tuned models get wrong. Their shared false positives follow
the vessel wall over an extended stretch and are systematically thinner than annotated plaques,
20.7 pixels against 29.9 in median thickness. Thickened wall without focal protrusion is exactly
the appearance that an enriched cohort supplies in quantity and a population cohort does not, and
it is exactly the case in which the distinction between plaque and not-plaque is hardest.

The asymmetry is not peculiar to this application. A lung ultrasound classifier developed at a
single centre and evaluated on multicentre data retained a sensitivity of 0.92 while its
specificity fell to 0.76, and fine-tuning on part of the external data raised specificity to 0.82
without changing sensitivity \cite{wu2024generalizability}. The correspondence with the figures
reported here is close: the base model retains 83.3\,\% sensitivity against 62.5\,\%
specificity, and fine-tuning recovers specificity while sensitivity does not decline. What
transfers across data sources appears to be the appearance of the lesion; what does not is the
threshold separating it from a diseased but lesion-free background, which is a property of the
population a model was trained on rather than of the pathology it was trained to recognise.

\subsection{What the Shift Would Do to a Prevalence Estimate}

The practical consequence for a phenotyping pipeline follows from the operating characteristics
alone. A classifier reports a positive rate of $\text{Se} \cdot p + (1 - \text{Sp})(1 - p)$ for a
true prevalence $p$. At the in-cohort values of 89.5\,\% and 89.2\,\%, a cohort with 45\,\% true
prevalence would be reported at 46.2\,\% -- an error of roughly one percentage point. At the
external values of 83.3\,\% and 62.5\,\%, the same cohort would be reported at 58.1\,\%.

The calculation is an illustration of the mechanism rather than a projection. It assumes that
sensitivity and specificity measured on a clinically enriched sample carry over to a population
sample, which spectrum effects make unlikely -- and the preceding subsection argues that precisely
this assumption fails here. What it does show is that a degradation of this size does not merely
make individual classifications less reliable; it biases any aggregate derived from them in a
predictable direction.

\subsection{The Shift Is Not One Thing}

Country, centre, device generation, operator, case mix and acquisition geometry differ
simultaneously between the two datasets, and a single external dataset cannot separate their
contributions. The heterogeneity documented in Section~\ref{subsec:external-dataset} shows that
the last of these varies appreciably even within Ar-PlaqSegm1, across groups whose median
grey-scale intensity differs by nearly a factor of two. What this thesis establishes is the
magnitude of the aggregate shift between two specific datasets, and that it is large enough to
matter for downstream use. Attributing it to particular causes would require external datasets
that vary one factor at a time, which do not currently exist for this task.

\section{What Limits the Fine-Tuned Models}
\label{sec:discussion-calibration}

Fine-tuning on target-domain data recovers roughly half the specificity lost to the shift, and
pooling two models recovers a little more. Neither addresses what actually constrains the models,
which the error analysis identifies as something other than segmentation quality.

\subsection{The Masks Are Already There}

Selecting, for each plaque-positive test image, the best-matching predicted instance irrespective
of its confidence raises the mean Dice coefficient from 0.655 to 0.794 and leaves no image
without a correct mask. The two images the pooled model appears to miss are both outlined
accurately -- one at a Dice of 0.82 -- and both carry confidence scores three orders of magnitude
below the operating threshold. Whatever limits these models, it is not their ability to delineate
a plaque once they have attended to it.

That this traces to a single training hyperparameter is worth stating plainly, because it is
unglamorous and consequential in equal measure. With one class, the classification branch of the
loss \emph{is} the confidence score, and it was weighted twenty-five times below the terms
governing geometry. Raising the weight moved the median confidence of the correct instance from
0.060 to 0.602. It also lowered the attainable Dice from 0.769 to 0.680 and left two images
without any correct mask, so the setting trades one failure mode for the other rather than
removing either.

\subsection{Ranking Is Not the Binding Constraint}

A natural response is to repair the ordering after training, and the attempt is instructive in
failing. A re-scorer using geometry, position and local echo contrast improved the separation of
correct from incorrect instances substantially, from an area under the curve of 0.681 to 0.823,
and recovered both previously missed images. Image-level performance nonetheless fell, from
$J = 0.667$ to 0.583, because specificity collapsed to 58.3\,\%.

The reason is that the two objectives are different. Features describing what a plaque looks like
identify genuine instances well, but a plaque-free image still elicits plaque-shaped proposals,
and the same features promote those equally. The original confidence score, for all its
miscalibration, encodes something the appearance features do not: whether there is anything there
at all. The binding constraint is objectness, and no reordering of proposals can supply it.

\subsection{What the Shared Errors Suggest}
\label{sec:discussion-shared-errors}

Five of the six false positives are produced by both models. Two networks differing in capacity,
input resolution and training schedule converge on the same mistakes, which points away from the
models and towards the data or the task definition. Under the Mannheim criterion \cite{touboul2012mannheim} a plaque is a
focal protrusion exceeding 50\,\% of the \emph{surrounding} intima-media thickness -- a relative
judgement requiring a comparison between a candidate region and its neighbourhood. A model that
scores each region on its own appearance has no access to that comparison, and thickened
non-focal wall is the case in which the omission shows.

One candidate explanation can be set aside. Annotation variability is not the limiting factor:
one image of the dataset occurs twice and carries two independently produced masks, which agree
at a Dice coefficient of 0.971 -- well above the 0.62 to 0.79 the models achieve
(Section~\ref{sec:results-duplicates}). A single pair is weak evidence, but it points away from label noise as the
ceiling and towards the models themselves.

If this reading is right, the shared error rate is a rough indication of what cannot be removed
by better training under the present label definition, and improvement would have to come from
encoding the criterion rather than from more capacity or more data. The inference is suggestive
rather than established: two models is a thin basis for a claim about irreducible error, and the
alternative explanation -- that both inherited the same bias from a shared COCO initialisation
and a common augmentation scheme -- is not excluded by the evidence here.

\subsection{A Methodological Note}

The re-scoring result generalises beyond this application. An intervention that improved the
ranking metric by 0.14 area under the curve degraded the decision that the ranking exists to
support. Calibration methods are conventionally assessed on calibration metrics; where the model
feeds a downstream decision, that decision is the only evaluation that settles whether the
intervention helped.

\section{Comparison with the Literature}
\label{sec:discussion-literature}

\subsection{How Much Non-Image Data Contributes}

\textcite{stamos2025} addresses a question adjacent to Phase 1A: whether clinical and biochemical
data, including protein biomarkers, add information to B-mode carotid ultrasound. Three fusion
architectures were evaluated on 73 patients stratified into high- and low-risk groups by degree
of stenosis. The result most relevant here is the decomposition. Imaging alone reached an area
under the curve of 84.6\,\%; the tabular data alone, 64.6\,\%; their fusion, 86.1\,\%. A learned
attention mechanism assigned 69.4\,\% of its weight to the imaging arm.

Two observations follow. First, non-image data carrying only modest standalone signal for a
carotid plaque phenotype is not peculiar to the present work: 64.6\,\% there and the 54.1\,\%
obtained from proteomics here are different tasks on different data, but they occupy the same
regime, well below what the image supplies. Second, and more consequentially, properly fused
end-to-end the tabular data added 1.5 percentage points to an imaging-only model. If that is the
scale of the available gain, then improving the imaging arm is the larger lever, which is the
direction this thesis took for unrelated reasons.

\subsection{Two Ways of Being Data-Limited}

The two studies fail in opposite directions, and the contrast is instructive. \textcite{stamos2025}
has clinically adjudicated risk labels and 73 patients: the ground truth is sound and the sample
is too small to support strong conclusions, as an area under the curve of 86.1\,\% on 73 cases
carries a wide interval. The present work has 20{,}014 participants and a phenotype produced by a
model whose own external validation shows substantial error. Neither configuration permits the
analysis both are reaching for. A study combining adjudicated labels with cohort-scale numbers
does not presently exist for carotid plaque, and until it does, results in this area will be
limited either by their labels or by their size.

It is also worth noting that \textcite{stamos2025} is an instance of the pattern
\textcite{omarov2024deep} identify as limiting the field: a cohort selected for carotid disease
rather than drawn from a population. The concern they raise about generalisability applies to it,
and symmetrically, the concern this thesis raises about phenotype quality applies to their own
population-scale results.

\subsection{Why the Obvious Comparison Is Unavailable}

The fine-tuned ensemble reported here attains 91.7\,\% sensitivity and 75.0\,\% specificity;
\textcite{omarov2024deep} report 89.5\,\% and 89.2\,\%. Setting these against one another would be
a mistake. They come from different datasets with different case mixes, acquisition conditions and
annotators, and the second set was obtained on the cohort the model was developed on while the
first was not. The only comparison this thesis licenses is the one made on a single test partition
in Section~\ref{sec:results-phase2}, where the base model and the fine-tuned models are evaluated
on identical images under an identical protocol.

The difficulty is not peculiar to these two studies. Every result cited in this chapter was
obtained on data held by the group that produced it: \textcite{omarov2024deep} on UK Biobank,
\textcite{stamos2025} on a cohort of 73 patients from one centre, and the present work on
Ar-PlaqSegm1. Each is internally coherent and none can be placed on a common scale with the
others. Reporting guidelines address the related but distinct question of whether a model was
validated outside its development cohort \cite{collins2015tripod, mongan2020claim}; they do
not make results from different external cohorts comparable to each other.

What distinguishes the evaluation reported here is that its test data are public.
Ar-PlaqSegm1 is openly available with a published partition \cite{rulloni2025arplaqsegm1}, so
every figure in Chapter~\ref{ch:results} can be recomputed by others and any future model can be
placed directly alongside them. That is a modest form of comparability, and extending it is the
subject of Section~\ref{sec:conclusion-outlook}.

\section{Limitations}
\label{sec:discussion-limitations}

\subsection{Adaptation, Not Formulation}
\label{subsec:limitations-control}

The comparison in Chapter~\ref{ch:results} places fine-tuned segmentation models against a
detection model that was never adapted to Ar-PlaqSegm1. It therefore confounds two things: the
change of formulation from bounding boxes to masks, and the adaptation to the target domain. A
control experiment was run to separate them. A detector was trained on the same images, the same
split and the same schedule as seg\_v2, differing only in its head and in the absence of
copy-paste augmentation, which requires masks and is unavailable to a box model.

The control reaches the same image-level performance as the segmentation models. At the threshold
selected on validation it attains a sensitivity of 79.2\,\% and a specificity of 83.3\,\%
($J = 0.625$) on the full test partition -- identical to seg\_v1, seg\_v2 and seg\_v3, with an
error profile coinciding exactly with that of seg\_v3. Excluding the duplicated images it reaches
$J = 0.583$, inside the range of 0.525 to 0.600 spanned by the four segmentation models. With 20
plaque-positive images one image corresponds to 0.05 in $J$, so every difference among the
fine-tuned models, the detector included, is smaller than a single image.

What this thesis demonstrates is consequently narrower than it might appear. The improvement over
the base model is attributable to adaptation to the target domain, not to the segmentation
formulation; a detector given the same data performs as well. The domain-shift result in
Section~\ref{sec:discussion-domain-shift} is unaffected, as is the calibration analysis, which
concerns the internal behaviour of the segmentation models rather than their advantage over
detection.

One argument for the segmentation formulation survives the control, because the control cannot
address it. A mask model yields lesion area and shape, and with them the Dice coefficients
reported throughout Chapter~\ref{ch:results}; a detector produces none of these. Where the
intended output is a binary phenotype the two formulations are interchangeable on this evidence,
and where it is a quantitative description of the lesion they are not.

A second control, a segmentation model trained without copy-paste augmentation, would have
separated the head from the augmentation within the segmentation formulation. It was started but
could not be completed, so the detector comparison above remains confounded with that
augmentation. Given that the detector matches the segmentation models despite lacking it, the
confound is unlikely to conceal an advantage rather than create one.

\subsection{Duplicate Images and a Contaminated Validation Split}

Two forms of leakage affect the results, one inherited and one introduced here.

The published split places four plaque-positive images in both partitions
(Section~\ref{subsec:duplicates}). Their effect on the reported figures is quantified in
Section~\ref{sec:results-duplicates} and is small, but it is leakage nonetheless, and any
published result computed on this test partition without deduplication carries the same
optimism.

The second is a consequence of how the development partition was divided here. The 493
development images were split into training and validation strata by plaque presence, but not by
image identity, and seven of the 22 duplicate groups internal to the development partition
straddle that division. Roughly one validation image in ten is therefore a copy of a training
image. Because the validation split selected both the best checkpoint and the confidence
threshold, both selections are mildly optimistic. Correcting this requires deduplicating before
splitting and retraining all configurations, which was not possible within the time available;
the appropriate response is to treat the validation figures as selection scores rather than
performance estimates, which is how they are used throughout.

\subsection{The Test Partition Does Not Span the Dataset}

The 48 test images are not a representative sample of Ar-PlaqSegm1. Grouping the images by
acquisition geometry (Section~\ref{subsec:external-dataset}) shows that 93.8\,\% of the test
partition falls into a single group, which makes up 52.1\,\% of the development partition, and
that three groups comprising 196 development images -- 37.7\,\% of the material available for
training -- are absent from the test partition altogether. The groups are visually and
statistically distinct, differing in median intensity by a factor of nearly two.

Every figure reported in Chapter~\ref{ch:results} therefore characterises performance on a narrow
region of the acquisition conditions the dataset contains. The models were trained across all
groups but evaluated on essentially one, and how they behave on the geometries absent from the
test partition is unknown. This qualifies the external validation in a direction that is easy to
overlook: evaluating on data from a different cohort establishes transfer across cohorts, but not
transfer across the acquisition conditions within the new cohort.

The observation also has a wider implication. If a published dataset that explicitly aims at
diversity -- its authors list the inclusion of varied sizes, locations and multiple plaques as a
design goal -- nonetheless distributes that diversity unevenly across its own partitions, then
comparable imbalances should be expected in evaluation sets generally, including the one on which
the base model was originally validated. Reporting the composition of a test partition, and not
only its size, would make such limitations visible.


\subsection{Statistical Resolution}

The held-out test partition contains 24 plaque-positive and 24 plaque-negative images, and after
excluding the duplicated images, 20 and 24. One image therefore corresponds to between 4.2 and
5.0 percentage points of sensitivity. Differences smaller than that are not measurable, and most
of the differences between the fine-tuned models are smaller than that: on the deduplicated set
the four segmentation models and the detection control span $J = 0.525$ to $0.600$, a range of
one and a half images.

Statements in this thesis are worded to respect that limit, and the conclusions that rest on
larger margins -- the gap to the base model, the size of the confidence-ranking effect -- are
unaffected. But any reading of the results that ranks the fine-tuned variants against one
another, or that treats the ensemble's advantage as established, exceeds what 44 images can
support. Cross-validation over the full development set would give a spread rather than a point
estimate, at a cost of a few training runs; it was not within the time available here.

\subsection{Patient-Level Independence}

The published split separates images, and the four duplicated images aside, no image occurs in
both partitions. The dataset record does not state how many patients the 541 image pairs derive
from, and patient-level independence has therefore not been established. If one patient
contributed images to both partitions, the models would have seen that patient's anatomy,
acquisition device and operator during training, and the test figures would be correspondingly
optimistic. The partition is the dataset authors' own rather than a choice made here, but the
limitation applies to the results regardless of who introduced it.

\subsection{Analytical Choices in Phase 1A}

Two decisions in the proteomic analysis would be made differently on repetition. Imputation and
standardisation were fitted on the full analysis cohort rather than within each cross-validation
fold, which leaks information across the folds; because both operations are unsupervised and the
bias runs towards optimistic performance, the null result reported is if anything conservative.
More consequentially, the clinical baseline used age at recruitment where age at the imaging
visit is the variable matched to the phenotype. The interval between the two is several years and
varies between participants, so the covariate is both shifted and noisy relative to the outcome
it is paired with, and the weak performance of the age-and-sex baseline cannot be cleanly
attributed to the phenotype until the analysis is repeated with the correct variable.

\subsection{One External Dataset}

External validation here rests on a single independent dataset. What the results establish is
that the base model transfers imperfectly to Ar-PlaqSegm1, and that fine-tuning on Ar-PlaqSegm1
recovers part of the loss. Whether the same would hold for a third dataset -- and how much of the
observed shift is attributable to country, centre, device, case mix or annotation practice --
cannot be determined from one comparison. A conclusion about external validity in general would
require several external datasets, as set out in Section~\ref{sec:conclusion-outlook}.

Annotation practice is the one component for which a partial estimate is available. The
duplicated image pair carrying two independently produced masks
(Section~\ref{sec:results-duplicates}) shows agreement at a Dice coefficient of 0.971, which
suggests that annotation variability within this dataset is small relative to the model errors
reported. A single pair is a thin basis for that, and no inter-reader agreement statistics are
published with the dataset.
